\documentclass[11pt]{article}

\usepackage[a4paper,margin=1.05in]{geometry}
\usepackage{mathpazo}
\usepackage[T1]{fontenc}
\usepackage{amsmath,amssymb}
\usepackage{microtype}
\usepackage{booktabs}
\usepackage{tabularx}
\usepackage{enumitem}
\usepackage{parskip}
\usepackage{titlesec}
\usepackage{fancyhdr}
\usepackage{tcolorbox}
\usepackage[colorlinks=true,linkcolor=accent,urlcolor=accent,citecolor=accent]{hyperref}

\tcbuselibrary{skins,breakable}

\definecolor{accent}{HTML}{1F5C6B}
\definecolor{accentbg}{HTML}{DCE9EC}
\definecolor{softgray}{HTML}{6B6B6B}
\definecolor{warnrule}{HTML}{B8860B}
\definecolor{warnbg}{HTML}{FBF0D9}

\titleformat{\section}{\large\bfseries\color{accent}}{\thesection}{0.7em}{}
\titleformat{\subsection}{\normalsize\bfseries}{\thesubsection}{0.6em}{}
\setlist[itemize]{leftmargin=1.3em,itemsep=2pt,topsep=3pt}
\setlist[enumerate]{leftmargin=1.6em,itemsep=2pt,topsep=3pt}

\pagestyle{fancy}
\fancyhf{}
\renewcommand{\headrulewidth}{0.3pt}
\lhead{\small\color{softgray}What the data is, and what model fits it}
\rhead{\small\color{softgray}\thepage}

\newtcolorbox{notebox}{
  enhanced, breakable, colback=accentbg!55, colframe=accent,
  boxrule=0pt, leftrule=2.5pt, arc=1pt, left=10pt, right=10pt, top=8pt, bottom=8pt}

\newtcolorbox{assumpbox}{
  enhanced, breakable, colback=warnbg!60, colframe=warnrule,
  boxrule=0pt, leftrule=2.5pt, arc=1pt, left=10pt, right=10pt, top=8pt, bottom=8pt}

\newcommand{\Ncal}{\mathcal{N}}

\title{\bfseries What the data is, and what model fits it}
\author{}
\date{}

\begin{document}
\maketitle
\vspace{-2.5em}

\begin{center}\small\color{softgray}
A sketch, not a specification. It fixes the shape of the data first, then writes
down the smallest model that consumes it.
\end{center}

\section{What a source gives you}

For one observable, in one cohort, a published paper gives one of two things:

\begin{itemize}
\item \textbf{the raw values}, where a figure plots per-patient points; or
\item \textbf{a few statistics of those values}: a median, a mean, an IQR, an SD,
      a range, a pair of quartiles.
\end{itemize}

Plus $n$ in both cases. Sometimes it also states an \emph{uncertainty} about one
of those statistics, such as a confidence interval on the median. Section
\ref{sec:distinction} explains why that is a different kind of object and not a
third statistic.

That is the entire inventory. Raw values if we are lucky, otherwise a handful of
summaries, and a sample size.

\begin{notebox}
\textbf{A statistic is a function, and that is the whole design.} Write $T$ for a
map from a sample to a number. The paper reports $T$ applied to its patients. The
model computes $T$ applied to its simulated cloud. The likelihood compares the
two.

Everything downstream follows from taking that literally. The record stores a
value and \emph{which} $T$. Nothing is converted, because nothing needs to be:
whatever the paper chose to report, the model can compute the same thing.
\end{notebox}

\subsection*{What we actually hold}

Measured over the 49 marginal targets currently in the set:

\begin{center}\small
\begin{tabularx}{\textwidth}{@{}Xl@{}}
\toprule
Quantity & Coverage \\
\midrule
$n$, a location, and exactly one further statistic & 49 / 49 \\
Per-patient values, so the shape is real & 21 / 49 \\
A scale that was extracted rather than generated from a summary & 21 / 49 \\
\bottomrule
\end{tabularx}
\end{center}

One scale statistic per target, never two. The \texttt{empirical\_data} block of
a target YAML holds a \texttt{median}, a single \texttt{ci95}, a
\texttt{sample\_size}, and, for 21 of the 49, the per-patient values behind them.

The other 28 targets are marked \texttt{synthesized}, which means their quantiles
were read off a distribution the pipeline built rather than off digitized
patients. Their \emph{shape} is our assumption. Their location and scale still
carry information, and sometimes real derivation work: \texttt{tam\_density} is
the sum of two separately reported cell populations, combined under a
marginalised correlation. What is not a measurement is the four quantile columns
that come out the far end.

The 21 targets with real per-patient values come from four cohorts:

\begin{center}\small
\begin{tabularx}{\textwidth}{@{}Xll@{}}
\toprule
Source & Observables & Patients \\
\midrule
Zheng 2022, GVAX + nivolumab arm & 8 & 10 \\
Zheng 2022, GVAX arm & 6 & 9 \\
Li 2022, baseline & 2 & 16 \\
Heumann 2023, urelumab arm & 1 & 8 \\
\bottomrule
\end{tabularx}
\end{center}

The same patients appear across the panels of a figure, but which point is which
patient is not recoverable from the figure. So these give real marginals and
give nothing about the joint.

\section{The one distinction that matters}
\label{sec:distinction}

Papers print two kinds of number and do not distinguish them:

\begin{itemize}
\item a \textbf{statistic}: something computed from the patients. A median, an
      SD, an IQR, a range, a pair of quartiles.
\item an \textbf{uncertainty about a statistic}: how well the study pinned one of
      the above. An SEM, a confidence interval on the median.
\end{itemize}

The second is not a third statistic. It is a statement about an estimator, and it
belongs in a different slot. A statistic goes on the left-hand side of a
likelihood row. An uncertainty does not go anywhere, because the model computes
its own from the cloud and $n_c$.

The two are related by $\sqrt n$, so they are easy to confuse and easy to
convert:
\begin{equation}
\mathrm{SE}(\text{median}) \;\approx\; 1.2533\,\frac{\sigma}{\sqrt n},
\qquad
\mathrm{SE}(\text{mean}) \;\approx\; \frac{\sigma}{\sqrt n}.
\label{eq:convert}
\end{equation}

\begin{notebox}
\textbf{What a reported uncertainty is good for.} One thing: checking. Compare it
against the bootstrap prediction from the statistic and $n$. Agreement validates
the extraction. Disagreement means the $n$ is wrong, the shape is unusual, or an
SEM was recorded as an SD.

It must not enter the likelihood. It carries no information the statistic and $n$
do not already carry, and feeding both counts one measurement twice.
\end{notebox}

\begin{assumpbox}
\textbf{One field, \texttt{ci95}, currently holds both kinds.} The distinction
above is not hypothetical here. The same field holds a statistic in some targets
and an uncertainty in others.

\emph{An uncertainty.} \texttt{cd8\_pct\_nonLA\_gvax\_d21} stores nine per-patient
values, which sort to $[0.0107,\,0.0207,\,0.0232,\,0.0795,\,0.0804,\,0.0838,\,
0.1022,\,0.1123,\,0.1493]$, and \texttt{ci95: [0.0207, 0.1123]}. Those are ranks
$2$ and $8$: the distribution-free order-statistic interval for a median, with
$96.1\%$ coverage at $n=9$. The same rank pattern appears at $n=7$. So
\texttt{ci95} there is a confidence interval on the location, computed from the
values stored beside it.

\emph{A statistic.} \texttt{tam\_density} stores \texttt{median: 110.44} and
\texttt{ci95: [18.64, 698.94]} at $n=113$. That is a ratio of $37.5$, an implied
log-scale sigma near $0.93$. A confidence interval on the median at $n=113$ would
have a ratio near $1.5$. So \texttt{ci95} there is a pair of patient-level
quantiles, which is two statistics.

The two readings differ by $\sqrt n$. The trial-arm targets, where the confusion
bites, run $n=6$ to $10$, so the factor is about $2.5$.

Under section \ref{sec:record} this cannot happen, because the two go in
different slots and each records which $T$ it is.
\end{assumpbox}

\section{The model that fits this}
\label{sec:model}

One row per reported statistic. Four steps.

\begin{enumerate}
\item \textbf{Simulate.} Draw $\vartheta_i\sim\Ncal(\mu,\Omega)$ with
      $\Omega=D_\omega R D_\omega$, push through the mechanism,
      $x_i=g(\vartheta_i)$. That is the predicted population of observables.
\item \textbf{Correct the predicted distribution.} The model can be wrong about
      where an observable sits and about how spread out it is, so give it one
      term for each:
      \begin{equation}
      \tilde x_{im} \;=\; \bar x_m \;+\; \kappa_m\,(x_{im}-\bar x_m) \;+\;
      \delta_m ,
      \qquad \bar x_m = \mathrm{med}_i\, x_{im}.
      \label{eq:correct}
      \end{equation}
\item \textbf{Apply the same $T$ the paper applied.} For each statistic $k$ that
      the source reports, compute $T_k$ on the corrected cloud.
\item \textbf{Compare.}
      \begin{equation}
      \widehat T_{mck} \;\sim\; \Ncal\!\Big(T_k\big(\tilde x_{\cdot m};\,\phi\big),
      \;\; V_{mck}\Big).
      \label{eq:row}
      \end{equation}
\end{enumerate}

\begin{center}\small
\begin{tabularx}{\textwidth}{@{}llX@{}}
\toprule
Term & Source & Meaning \\
\midrule
$\widehat T_{mck}$ & data & The value the paper reports for statistic $k$. \\
$T_k$ & data & \emph{Which} statistic that is. A function of a sample. \\
$n_c$ & data & The cohort size. \\
$\delta_m$ & inferred & Discrepancy of the model in the location of observable $m$. \\
$\log\kappa_m$ & inferred & Discrepancy of the model in the scale of observable $m$. \\
$V_{mck}$ & derived & Sampling variance of $T_k$ at $n_c$. Section \ref{sec:noise}. \\
\bottomrule
\end{tabularx}
\end{center}

Note where \eqref{eq:correct} sits. The discrepancy acts on the predicted
\emph{distribution}, before any statistic is taken, so it means something precise:
the model's predicted spread of observable $m$ is wrong by a factor $\kappa_m$
and its centre is wrong by $\delta_m$. Every statistic then inherits the
correction automatically, including one nobody has written yet. Two discrepancy
terms per observable, however many statistics that observable reports.

\section{Noise, and why to bootstrap it}
\label{sec:noise}

$V_{mck}$ is the sampling variance of $T_k$ computed from $n_c$ patients. Two
ways to get it.

\textbf{Analytic}, where a formula exists. The sampling variance of a log
standard deviation is about $1/\{2(n-1)\}$, whatever the population spread is. So
that row's noise is a known function of $n$ alone, free of $\phi$, which is worth
a lot: it removes the $-\tfrac12\log\det V$ term that would otherwise reward a
narrow population regardless of fit quality. Formulas of this kind exist for the
mean, the median and the SD, and for almost nothing else.

\textbf{Bootstrap}, always. Resample $n_c$ patients from the corrected cloud,
recompute $T_k$, repeat. That is exact at any $n$, assumes no shape, needs no
derivation, and works for any $T_k$ whatever.

\begin{notebox}
\textbf{The bootstrap is what makes custom statistics affordable.} Without it,
every new kind of reported statistic needs its own variance formula, and most do
not have one. With it, a statistic is defined once as a function and its noise
follows for free.

The cost is that a bootstrapped $V$ depends on $\phi$ through the cloud, which
brings back the $\log\det$ pathology. The usual remedy applies: evaluate it once
at a plug-in value and freeze it.

Suggested default: bootstrap uniformly and freeze, because uniformity is worth
more than the $\phi$-free property, and check it against the analytic forms where
those exist. The targets with $n=6$ to $10$ are where the two will disagree if
anywhere, because that is where the asymptotic formulas are least defensible.
\end{notebox}

\textbf{Rows are close to independent.} For a normal sample the mean and the
standard deviation are exactly independent, and on the log scale for a roughly
lognormal observable that holds well enough. Where two statistics of the same
sample are correlated, the bootstrap gives the covariance at no extra cost,
because it resamples the sample once and computes both.

\begin{notebox}
\textbf{Which statistic to ask for, when we get to choose.} More statistics from
one sample do carry more information, up to the limit of knowing all $n$ values,
and a correct $V$ makes sure each row contributes what it earned rather than
being counted as an independent measurement. But the \emph{choice} of statistic
matters far more than the number.

For a normal sample the SD is fully efficient for $\sigma$, with asymptotic
variance $\sigma^2/2n$. The interquartile estimator
$\widehat\sigma=\mathrm{IQR}/1.349$ has asymptotic variance $1.36\,\sigma^2/n$,
an efficiency near $37\%$. Reporting an IQR instead of an SD discards about two
thirds of what the sample says about the spread. At $n=10$ that is like working
with fewer than four patients.

So for the 21 targets that carry raw values, compute the log-scale standard
deviation, not the interquartile range. It costs nothing and it is worth more
than any amount of additional quantile extraction. Fall back to the IQR only
where the log-scale shape is skewed or heavy-tailed enough that robustness beats
efficiency, which is a thing the raw values let us check.

For the summary-only targets there is no choice. We match whatever was printed.
\end{notebox}

\begin{notebox}
\textbf{What this removes.} An earlier form fixed three quantile rows per
observable, at $0.25$, $0.50$ and $0.75$, whatever the paper had reported. Three
quantiles from one sample are strongly correlated, so it needed an asymptotic
quantile covariance, a density estimate at each quantile, a bandwidth, and a
Gaussian copula across observables. It also needed every source that did not
report quartiles to have quartiles manufactured for it.

None of that was solving a problem in the data. It was solving a problem created
by demanding a fixed statistic. Match whatever the paper reported and the
is diagonal.
\end{notebox}

\section{What is left}

$\mu$, $\omega$, and two discrepancy terms per observable with their scales
$\tau_\delta$ and $\tau_\kappa$. That is the model.

The name, so the literature is findable: this is \textbf{aggregate-data, or
summary-level, meta-analysis on a nonlinear forward model}. Equivalently it is a
nonlinear mixed-effects model fitted to published summaries instead of to
individual patient data. Pharmacometrics calls the same object model-based
meta-analysis.

\section{Deliberately left out of the first draft}

\textbf{Correlation between observables measured in the same patients.} It is
real: Zheng 2022 reports eight observables in ten patients. Ignoring it makes the
fit overconfident, because those rows are not independent evidence. It is a
refinement, not a foundation, and the linkage needed to estimate it is not in
hand.

\textbf{A study offset $\eta_c$.} Every observable in the set belongs to one
cohort, so a free per-study offset is exactly aliased with $\delta_m$. The simple
version omits it and lets $\delta_m$ absorb the study effect. Nothing is lost
that was identified.

\textbf{Fitting the shape itself.} Twenty-one targets have per-patient values,
so one could compare whole distributions rather than a few statistics. Resist for
now: the rows stop being independent, and the other 28 targets cannot be treated
the same way. Use the per-patient values to compute whichever statistics you
want, and to check whether the predicted shape is right.

\section{What the record has to hold}
\label{sec:record}

Follow the design of section \ref{sec:model} and the schema falls out. Per
observable, per cohort:

\begin{center}\small
\begin{tabularx}{\textwidth}{@{}lX@{}}
\toprule
Slot & Contents \\
\midrule
\texttt{n} & The number of patients. \\
\texttt{values} & The raw per-patient values, when a figure gives them. Best case:
                  every statistic below is then computable and checkable. \\
\texttt{statistics} & A list of (which $T$, value). \texttt{median},
                      \texttt{mean}, \texttt{sd}, \texttt{iqr}, \texttt{q25},
                      \texttt{q75}, \texttt{min}, \texttt{max} cover most of it. \\
\texttt{custom} & A statistic that the enum cannot name, written as a function of
                  a sample. \\
\texttt{uncertainty} & A reported SE or CI, tagged with which statistic it refers
                       to. Not a likelihood input. Used to check the bootstrap. \\
\bottomrule
\end{tabularx}
\end{center}

\begin{notebox}
\textbf{Why the custom slot is not optional.} An enum cannot express a compound
statistic without ambiguity, and the current set already contains several.
\texttt{m1\_m2\_ratio} is either the ratio of the medians or the median of the
ratios, and those are different numbers. \texttt{cd8\_fc} is a paired per-patient
fold change, so it is a statistic of two samples at once. A proportion of patients
above a threshold is a statistic no location-and-scale vocabulary can name.

Writing the function removes the ambiguity, and the same function then runs on
the simulated cloud. That is the contract: \emph{one function, two samples}. It
is testable, which a naming convention is not.
\end{notebox}

The \texttt{uncertainty} slot is the one that fixes the \texttt{ci95} problem in
section \ref{sec:distinction}. A patient-level 95\% interval is two entries in
\texttt{statistics}, namely \texttt{q2.5} and \texttt{q97.5}. A confidence
interval on the median is one entry in \texttt{uncertainty}, tagged
\texttt{median}. Today both are written to one field named \texttt{ci95} and
nothing distinguishes them.

\section{Checklist: cleaning up the data record}

One principle. \textbf{The record holds what the paper reported. Every
transformation happens at load, in one line, and is reversible.} Each item below
is a place where that is not true today.

\begin{itemize}[label={},leftmargin=1.6em]

\item[$\square$] \textbf{Stop materialising derived quantiles. Keep
\texttt{distribution\_code}.} The field does two different jobs, and only one is
cruft.

Sometimes it is a real derivation. \texttt{tam\_density} is the sum of the
CD163$+$ and CD163$-$ densities that Liudahl 2021 reports separately, each with a
median and quartiles. A sum of two lognormals has no closed form, so the code
builds both marginals, marginalises an unidentified within-patient correlation
$\rho\sim\mathrm{Beta}(a,b)$ through a one-factor copula, and samples. That is the
honest record of how a model observable was assembled from what the paper
printed, and it handles the unknown correlation by integrating it out rather than
assuming it. Keep it.

Sometimes it is format padding: a source reports a location and a scale, and the
code exists only to manufacture quantiles at $0.25$ and $0.75$. That output lands
in columns indistinguishable from digitized data, so nothing downstream can tell
which is which. The assumption is not the problem. Hiding it in the record is.

The fix is small. \texttt{distribution\_code} already ends in
\texttt{pop.summarize(samples)}, and a set of named statistics is exactly what
section \ref{sec:record} wants. Have it return those, and stop reading a fixed
quantile grid off the samples.

\item[$\square$] \textbf{Populate \texttt{dispersion\_type} on every scale
statistic.} The field exists on each input and is \texttt{null} in every target
checked. It needs one of two values, \texttt{dispersion} or \texttt{precision},
because \texttt{ci95} currently holds both and they differ by $\sqrt n$. For the
21 targets with per-patient values it can be set by inspection, since the raw
data sits in the same file.

\item[$\square$] \textbf{Drop the manufactured columns from the marginal-targets
CSV.} \texttt{q10}, \texttt{q25}, \texttt{q75}, \texttt{q90} and \texttt{iqr}
exist to serve a fixed three-quantile input format that the model in section
\ref{sec:model} does not ask for. Replace them with the slots of section
\ref{sec:record}. Keep \texttt{kind}, or something like it, as provenance: it is
the only marker of whether the shape came from patients or from us.

\item[$\square$] \textbf{Keep the per-patient values wherever they exist.} They
are the only check on the shape assumption, and the only route to anything about
the joint distribution across observables. Storing them costs nothing.

\end{itemize}

Model changes are a separate list. Nothing above depends on settling them.

\end{document}
